\subsection{مبانی نظری و فرمول‌بندی ریاضی}
در گذر پیشرو (\lr{Forward Pass})، توزیع توجه از طریق فرمولاسیون پایدار بیشینه‌ساز نرم بر روی بعد آخر ماتریس ضرب شباهت محاسبه می‌شود:
\begin{equation}
y_i = \frac{\exp(x_i - \max_k x_k)}{\sum_j \exp(x_j - \max_k x_k)}, \quad \forall i
\end{equation}
در گذر پسرو (\lr{Backward Pass})، مشتق تحلیلی نسبت به ورودی $x$ با استفاده از بردار گرادیان پس‌انتشاریافته $g_i = \frac{\partial \mathcal{L}}{\partial y_i}$ به فرم زیر تدوین می‌گردد:
\begin{equation}
\frac{\partial \mathcal{L}}{\partial x_i} = y_i \left( g_i - \sum_j g_j y_j \right) = y_i (g_i - \mathbf{g} \cdot \mathbf{y})
\end{equation}
به منظور جلوگیری از فروپاشی دقت در محاسبات نیم‌دقت (\lr{FP16/BF16})، مقادیر خروجی به صورت برخط مورد ارزیابی قرار گرفته و مقادیر غیرمعتبر (\lr{NaN/Inf}) پاکسازی شده و تحت بازه امن مهار می‌گردند:
\begin{equation}
\operatorname{Clamp}\left(\frac{\partial \mathcal{L}}{\partial x_i}\right) = \max\left(-C, \min\left(C, \frac{\partial \mathcal{L}}{\partial x_i}\right)\right), \quad C = 65000.0
\end{equation}
این ساختار مانع از انفجار گرادیان در توالی‌های طولانی با امتیازات شباهت بسیار بالا یا پایین می‌شود.
